\documentclass[a4paper]{article}

\usepackage[english]{babel}
%\usepackage[utf8x]{inputenc}
\usepackage{amsmath}
\usepackage{graphicx}
\usepackage[colorinlistoftodos]{todonotes}
\usepackage[all]{xy}
\usepackage{dsfont}
\usepackage{amssymb}
\usepackage{bbold}
\usepackage{amsthm}

\usepackage{tikz}
\usepackage{verbatim}
\usetikzlibrary{matrix}
\usetikzlibrary{cd}

\usepackage[autostyle=true]{csquotes}
\usepackage{pdfpages}
\usepackage{hyperref}
\usepackage{fancyhdr}



\newtheorem{theorem}{Theorem}[section]
\newtheorem{proposition}{Proposition}[theorem]
\newtheorem{lemma}{Lemma}[theorem]
\newtheorem{corollary}{Corollary}[theorem]
\newtheorem*{definition}{Definition}


\newcommand{\norm}[1]{\left\lVert#1\right\rVert}

\newcommand{\Rn}{\mathbb{R}^n}


\newcommand{\too}[1]{\overset{#1}{\longrightarrow}}
\newcommand{\hookar}[1]{\overset{#1}{\hookrightarrow}}
\newcommand{\colim}[1]{ \underset{#1}{\texttt{colim}} \text{ } }
\newcommand{\fset}[1]{ \{ #1_0 , \dots , #1_n \} }
\newcommand{\ddt}{\frac{d}{dt}}


\newcommand{\N}{\mathbb{N}}
\newcommand{\Z}{\mathbb{Z}}







\fancyhead[L]{\includegraphics[scale=0.1]{UFT.png}}

\begin{document}


\thispagestyle{fancy}
\vspace*{1cm}

\begin{minipage}[r]{12cm}
    \begin{flushright}
\begin{tikzpicture}
\draw[green, very thick] (0,0) rectangle (3,2);
\end{tikzpicture}
    \end{flushright}
\end{minipage}

\begin{center}
\Large{\textbf{Control Recuperativo \\\'Algebra Lineal\\  }}
\end{center}
\vspace{1cm}
\begin{minipage}[r]{15cm}
\begin{flushleft}

Fecha: 30 de noviembre de 2023 \\%
Duración: 120 minutos  \\
Prof: C. Muñoz.\\

\end{flushleft}
\end{minipage}
\vspace{2cm}
\begin{center}
\Large{\textbf{Facultad de Ingeniería}}
\end{center}
\vspace{1cm}
\textbf{En caso de que un alumno sea sorprendido en actitudes deshonestas, ya sea intentando o
dejándose copiar, será calificado con nota mínima 1.0 y perderá toda posibilidad de eliminar
dicha nota. Además, esta conducta quedará archivada como antecedente en su carpeta y
podrá ser causal de eliminación.
Así mismo, está prohibido salir de la sala, usar o mantener sobre la mesa teléfonos celulares
o cualquier tipo de aparato electrónico, a no ser que el profesor lo requiera.}\\



\vspace{3cm}
\makebox[12cm]{\textbf{Nombre Alumno}:\ \hrulefill} 
\newpage




\newpage
\thispagestyle{empty}
\begin{document}

\section*{Instrucciones}
\begin{itemize}
\item No se asignar\'a puntaje a los ejercicios sin desarrollo.
    \item Tiene 120 minutos para responder el control.
    \item No está permitido el uso de celular ni calculadora.
    \item No está permitido el préstamo de materiales.
\end{itemize}
\section*{Preguntas}





\begin{enumerate}
    \item Considere el conjunto $S=\{ \begin{pmatrix} 0 \\ 3 \\ 2 \end{pmatrix}, \begin{pmatrix} 1 \\ 1 \\ 2 \end{pmatrix}, \begin{pmatrix} 0 \\ 0 \\ 5 \end{pmatrix} \}$ 
    \begin{enumerate}
        \item Muestre que $S$ es base de $\mathbb{R}^{3}$.
        \item Encuentre una base ortonormal de $\mathbb{R}^{3}$ a partir de esta base usando el procedimiento de Gram-Schmidt.
        \item Encuentre las coordenadas de $\{ \begin{pmatrix} 1 \\ 1 \\ 1 \end{pmatrix} \}$ en esta base.
    \end{enumerate} 
\newpage    
    \item Encuentre la proyección de $\{ \begin{pmatrix} 1 \\ 1 \\ 1 \end{pmatrix} \}$ sobre el plano generado por $\{ \begin{pmatrix} 0 \\ 3 \\ 2 \end{pmatrix}, \begin{pmatrix} 1 \\ 1 \\ 2 \end{pmatrix} \}$.

\newpage
             \item Muestre que la función $T:\mathbb{R}^2\to \mathbb{P}_2[x]$ definida por $$T(a,b)=(2a+b)+(a-b)x+(a+b)x^2,$$ es una aplicación lineal. Luego determine $\ker T$ y $\dim \operatorname{Im} T$.
\newpage     
\item  Demuestre que esta función $T: \mathbb{R}^{3} \to \mathbb{R}^{3}$ dada por $T \begin{pmatrix} x \\ y \\ z \end{pmatrix} = \begin{pmatrix}
    2x+z \\ x-y+z \\ y+2z
\end{pmatrix}$
 es lineal y encuentre su matriz representante en la base canónica.
 \end{enumerate}
\newpage
\section*{Formulario} 
\begin{itemize} 
\item \textbf{Bases:} Un conjunto de vectores $\{ v_1, \dots , v_n \}$ se dice linealmente independiente si $ \lambda_1 v_1+ \dots + \lambda_n v_n= \Vec{0} \Longrightarrow \lambda_1= \dots = \lambda_n= 0$. $W$ se dice generado por $\{w_1, \dots , w_n\}$ si $<\{w_1, \dots , w_n\}>=\{ \sum \lambda_iw_i | \lambda_i \in \mathbb{R} \}=W$ Una base de un espacio vectorial $V$ es un conjunto linealmente independiente que genera $V$. \item \textbf{Procedimiento para obtener una base ortonormal:} Sea $B = \{v_i\}_{i=1}^m $ una base ordenada cualquiera de $W $. \begin{itemize} \item Escribimos $w_1 = v_1$. \item Calculamos $w_2=v_2-\dfrac{<v_2,w_1>}{<w_1,w_1>}w_1$. \item Calculamos $w_3=v_3-\dfrac{<v_3,w_1>}{<w_1,w_1>}v_1- \dfrac{<v_3,w_2>}{<w_2,w_2>}w_2$. \item Inductivamente obtenemos los vectores $w_m=v_m-\sum_{i=1}^{m-1}c_iw_i$ con $c_i=\frac{<v_m,w_i>}{<w_i,w_i>}$. \end{itemize} El conjunto $\{ w_1,w_2, \dots, w_n \}$ es una base ortogonal. Si se normaliza cada vector $w_k$, obtenemos nuevos vectores $u_k=\frac{w_k}{\|w_k\|}$. Luego, $\{ u_1,u_2, \dots, u_n \}$ una base ortonormal para $W$. 
\item \textbf{ Coordenada respecto de una base ortogonal } $B = \{w_i\}_{i=1}^m $ est\'a dada por $a_i=\dfrac{\langle v, v_i \rangle}{\langle v_i, v_i \rangle}$ \item {\bf Proyección ortogonal: } Sea $W$ un subespacio de $\mathbb{R}^n$ con base ortogonal $\{ w_1,w_2, \dots , w_k \}$. Si $v \in \mathbb{R}^n$, entonces la proyección ortogonal de $v$ sobre $W$, denotada por $P_W(v)$, está dada por: $$P_W( v) = \frac{\langle v, w_1 \rangle}{\langle w_1, w_1 \rangle}w_1 + \frac{\langle v, w_2 \rangle}{\langle w_2, w_2 \rangle}w_2 + \cdots + \frac{\langle v, w_k \rangle}{\langle w_k, w_k \rangle}w_k.$$ \item {\bf Transformaciones lineales: } Transformación Lineal $T$ de $ V$ en $W$, como una función $T:V\to W$ tal que \begin{enumerate} \item $T(v+w)=T(v)+T(w)$, \item $T(\alpha v)=\alpha T(v)$, \end{enumerate} para todo $v,w\in V$, para todo $\alpha \in K$. \item $T:V\to W$ transformaci\'on lineal, entonces $$\operatorname{dim}V=\operatorname{dim} \operatorname{ker }T+\operatorname{dim}\operatorname{Im}T$$ \item {\bf Valores y Vectores Propios: } $\lambda \in \mathbb{R}$ es valor propio de $M \in \mathcal{M}_{n \times n}$ si existe $\Vec{v}\in \mathbb{R}^n$ con $M\Vec{v}= \lambda \Vec{v}$. $\Vec{v}$ se llama vector propio asociado a $\lambda$ \begin{itemize} \item Los valores propios son las raices del polinomio característico $p(x)=\texttt{det}(M - x I_n ) $ \item Los vectores propios asociados a un valor propio $ \lambda_0$ forman el subespacio propio $E_{\lambda}=\texttt{Ker}(M-\lambda_0 I_n)$ \item La multiplicidad geométrica de $\lambda$ es $\texttt{dim}(E_\lambda)$ y la multiplicidad algebráica es su multiplicidad como raiz del polinomio característico. \end{itemize} \end{itemize}

\end{document}